% Part: computability
% Chapter: computability-theory
% Section: ce-sets

\documentclass[../../../include/open-logic-section]{subfiles}

\begin{document}

% SEGMENT OLP-0238-S01

\olfileid{cmp}{thy}{ces}
\olsection{கணிக்கத்தக்கவாறு எண்ணிடத்தக்க கணங்கள்}

\begin{defn}
ஒரு கணம் வெற்றுக் கணமாகவோ ஒரு கணிக்கத்தக்க சார்பின் வீச்சாகவோ இருந்தால்
அது \emph{கணிக்கத்தக்கவாறு எண்ணிடத்தக்கது}.
\end{defn}

\begin{history}
கணிக்கத்தக்கவாறு எண்ணிடத்தக்க கணங்கள்
\emph{மீள்வரையறை முறையில் எண்ணிடத்தக்கவை} என்றும் அழைக்கப்படுகின்றன.
இதுவே மூலக் கலைச்சொல்; இன்று இரு பெயர்களும், “c.e.” மற்றும் “r.e.”
என்ற சுருக்கங்களும் பொதுவாகப் பயன்படுத்தப்படுகின்றன.
\end{history}

\begin{explain}
இவ்வரையறை எதைக் குறிக்கிறது என்றும், கலைச்சொல் ஏன் பொருத்தமானது
என்றும் சிந்திக்க வேண்டும். $S$ என்பது கணிக்கத்தக்க
சார்பு~$f$ இன் வீச்சு எனில்,
\[
S = \{ f(0), f(1), f(2), \dots \},
\]
எனவே~$S$ இன் உறுப்புகளை $f$ “எண்ணிடுவதாகக்” கருதலாம் என்பதே
கருத்து. வரையறைப்படி, $f$ ஓர் ஏறு சார்பாக இருக்கத் தேவையில்லை;
அதாவது எண்ணிடல் ஏறு வரிசையில் இருக்கத் தேவையில்லை என்பதைக் கவனிக்க.
உண்மையில் $f$ ஓருறுப்புச் சார்பாகவும் இருக்கத் தேவையில்லை;
அதாவது~$S$ இன் எண்ணிடல் $f(0)$, $f(1)$, $f(2)$, \dots{}
இல் மீள்தோற்றங்கள் அனுமதிக்கப்படுகின்றன. எடுத்துக்காட்டாக,
மாறிலிச் சார்பு $f(x) = 0$ ஆனது கணம் $\{ 0 \}$ ஐ எண்ணிடுகிறது.
\end{explain}

ஒவ்வொரு கணிக்கத்தக்க கணமும் கணிக்கத்தக்கவாறு எண்ணிடத்தக்கது.
இதைக் காண,~$S$ கணிக்கத்தக்கது என்று கொள்க. $S$ வெற்றுக் கணம்
எனில், வரையறைப்படி அது கணிக்கத்தக்கவாறு எண்ணிடத்தக்கது.
இல்லையெனில்,~$S$ இன் ஏதேனும் உறுப்பு $a$ ஐக் கொள்க. $f$ ஐப்
பின்வருமாறு வரையறுக்க:
\[
f(x) =
\begin{cases}
x & \text{$\Char{S}(x) = 1$ எனில்} \\
a & \text{இல்லையெனில்.}
\end{cases}
\]
அப்போது $f$ ஒரு கணிக்கத்தக்க சார்பு; மேலும்~$S$ ஆனது~$f$ இன் வீச்சு.

\end{document}
